\begin{lemma}
Let $(\Omega,\nu)$ be a probability space, $A,K$ finite types, and give
$A$ a sigma algebra with measurable singletons. Let $X_a:\Omega\to A$
be measurable independent random variables. Let $P\subseteq A$ satisfy
$X_a(\omega)\in P$ for every $\omega,a$. Suppose
$f:A^K\to\mathbb R$ satisfies $|f(I)-f(J)|\le1$ whenever $I,J\in P^K$
agree outside one coordinate. If $m=|K|>0$ and $t>0$, then
\[
\nu\{f(X)\ge\mathbb E f(X)+t\}\le \exp(-2t^2/m).
\]
\end{lemma}
\begin{proof}
All subsets of the finite set $A$ are measurable. Write
$p_a(v)=\nu\{X_a=v\}$. Independence gives
$\nu\{X=I\}=\prod_a p_a(I_a)$, and each $p_a$ sums to one.
Thus the push-forward of $\nu$ by $\omega\mapsto(X_a(\omega))_a$
is the finite product of these masses. Explicitly, for every $B\subseteq A^K$
and every real function $g$ on $A^K$,
\[
\nu\{X\in B\}=\sum_{I\in B}\prod_a p_a(I_a),\qquad
\int g(X)\,d\nu=\sum_{I\in A^K}g(I)\prod_a p_a(I_a).
\]
These identities follow by partitioning into the finitely many measurable
fibers; the integral is legitimate because $g(X)$ has finite range.
Every value with positive $p_a$ belongs to $P$, since its fiber is
nonempty. We may therefore restrict each coordinate to its nonempty
positive-mass support without changing either identity.

Enumerate $K$ by $1,\ldots,m$. For a supported prefix $v_1,\ldots,v_i$,
define $M_i$ to be the sum of $f(v_1,\ldots,v_i,w_{i+1},\ldots,w_m)$
against the product of the remaining coordinate masses. Thus $M_0=\mathbb Ef(X)$,
$M_m=f(X)$, and averaging $M_i$ over $v_i$ gives $M_{i-1}$.
For a fixed prefix of length $i-1$, two possible values of $M_i$ differ
by at most one: subtract their defining sums, use the assumed pointwise
bound on $f$, and use that the remaining product masses sum to one.
If $u$ is the minimum of these values, all lie in $[u,u+1]$.
Consequently $D_i=M_i-M_{i-1}$ has conditional mean zero and lies
in $[-p,1-p]$, where $p=M_{i-1}-u\in[0,1]$.

Here is the needed interval-width-one exponential estimate. Convexity
gives, for $z\in[-p,1-p]$ and $\lambda\ge0$,
\[
e^{\lambda z}\le(1-p-z)e^{-\lambda p}+(z+p)e^{\lambda(1-p)}.
\]
Averaging a mean-zero $z$ bounds its exponential mean by
$B(\lambda)=(1-p)e^{-\lambda p}+pe^{\lambda(1-p)}$.
For $p=0$ or $1$ this is one. Otherwise $h=\log B$ has
$h(0)=h'(0)=0$ and $h''(\lambda)=q(1-q)\le1/4$, where
$q=pe^{\lambda(1-p)}/B(\lambda)$. Integrating this inequality twice
on $[0,\lambda]$ gives $h(\lambda)\le\lambda^2/8$.
Thus every prefix has conditional exponential mean at most
$e^{\lambda^2/8}$. Successively summing over the last exposed coordinate
in the finite product shows
\[
\mathbb E e^{\lambda\sum_iD_i}\le e^{m\lambda^2/8}.
\]
For $\lambda>0$, on the non-strict tail event the integrand is at least
$e^{\lambda t}$, so its probability is at most
$e^{-\lambda t+m\lambda^2/8}$. Choose $\lambda=4t/m>0$ and use
$\sum_iD_i=f(X)-\mathbb Ef(X)$ to obtain the stated bound. The fiber
identities above transport both the mean and the event back to $\Omega$.
\end{proof}
